\answersfor{ch:hslunmix}

\begin{solution}{ex:hslunmix:1}
At the target, $(0.25+8\times0)/9=0.25/9=0.0278$. A diagonal neighbour's window also contains the target once, so it is
0.0278 too. A pixel two steps away has a window that does not reach the target: 0.
\end{solution}

\begin{solution}{ex:hslunmix:2}
\begin{align*}
\text{upstream: } 0.106\times57.296 &= 6.07^\circ,\\
\text{VCA + FCLS: } 0.004\times57.296 &= 0.23^\circ .
\end{align*}
\end{solution}

\begin{solution}{ex:hslunmix:3}
The target survives while $0.13-2\lambda/0.20>0$, that is $2\lambda<0.13\times0.20=0.026$, so $\lambda<0.013$. The
$\lambda=0.013$ of the worked example is exactly the boundary, which is why the target landed on 0.
\end{solution}

\begin{solution}{ex:hslunmix:4}
For $k=2$ the shrinkage is $2\lambda/(kc)=0.026/(2\times0.20)=0.026/0.40=0.065$, so $a'=0.13-0.065=0.065$. It survives, at
half its fill: four pixels share a perimeter of 8 edges, half the perimeter per pixel of a lone pixel.
\end{solution}

\begin{solution}{ex:hslunmix:5}
If the two parts are at right angles to each other\tagA, squared lengths add, and $\mat{E}\vect{a}$ is left with
$1.0-0.40=0.60$. The abundances need not account for the whole pixel: whatever the unconstrained branch can absorb,
including a target's signature, no longer has to show up as a fill.
\end{solution}

\begin{solution}{ex:hslunmix:6}
$4\times13=52$ minutes for the bench; $18\times13=234$ minutes, 3.9 hours, for the implant hosts, against seconds per scene
for VCA + FCLS.
\end{solution}
